\scanpage{170}
\mspage{86}

Nous allons reformuler les données (\unsure{39|34})(36) et (41) (le scindage \deleted{\unsure{$g$|$\xi$}} de l'extension $B$ de $Z$ par $B'$) en \deleted{\unsure{inter}}\inserted{les} réinterprétant \deleted{ce diagramme comme \unsure{donnée}} \inserted{comme l'une des} \deleted{\unsure{don}}nées suivantes

\noindent\begin{tabular}{@{}l@{}l@{}}
(42)\ $\left\{\vphantom{\rule{0pt}{11cm}}\right.$ & \begin{minipage}{.84\linewidth}
a) Un diagramme \inserted{carré} de groupes, cartésien et cocartésien, avec des inclusions de groupes qui sont \unsure{distingués|distingué}

\diagram{\begin{tikzcd}[ampersand replacement=\&,row sep=small]
\widetilde B''\arrow[r,hook,"Z'"]\&B\\
\widetilde U\arrow[u,hook,no head,"Z"]\arrow[r,hook,"Z'"]\&B'\arrow[u,hook,"Z"']
\end{tikzcd}}

b) Un scindage de l'extension $B$ de $Z$ par $B'$, avec un groupe-section $Z\subset B$ central dans $B$, de sorte que l'on a un iso. d'ext.

\[ B\simeq B'\times Z,\qquad \deleted{\unsure{B}}Z\simeq\widetilde B''/\widetilde U\ \text{commutatif} \]

c) Un \deleted{opération de} \inserted{plongement} de groupes

\[ \widetilde B''\mathrel{\unsure{\lhook\joinrel-|\hookrightarrow}}G_0 \]

d) Une opération de $B$ sur $G_0$,

\[ B\longrightarrow\operatorname{Aut}_{\rm gr}(G_0) \]

prolongeant l'opération tautologique de $\widetilde B''$ sur $G_0$ déduite de c), et \uline{triviale} sur le sous-groupe $Z$ de $B$. (opération qui permet alors de réconstituer le groupe $G$ et le diagramme (39))
\end{minipage}\end{tabular}

Nous allons encore modifier un peu cette traduction, en écrivant le diagramme (42 a) \inserted{$+$ donnée b)} sous la forme

\diagram{\begin{tikzcd}[ampersand replacement=\&,row sep=small]
(43)\quad\&\widetilde B''\arrow[r,hook,"{\deleted{\unsure{Z'}}}","{\unsure{i|t}}"']\&B'\times Z\&\\
\&\widetilde U\arrow[u,hook,"Z"]\arrow[r,hook,"Z'"]\&B'=B'\times\{0\}\arrow[u,hook]\& Z=\widetilde B''/\widetilde U
\end{tikzcd}}

% archive sheet number: 169
